\begin{resumo}
Este trabalho propõe investigar a relação entre exposição ocupacional à inteligência artificial e transições no mercado de trabalho brasileiro, utilizando a PNAD Contínua. A pesquisa considera mudanças de ocupação, saída do emprego e passagem para a informalidade, com o lançamento público do ChatGPT como referência temporal para a comparação entre períodos. O desenho utiliza pares de entrevistas consecutivas, medidas de exposição ocupacional e características dos trabalhadores. Os procedimentos compreendem a preparação dos insumos, a correspondência entre classificações ocupacionais, a construção do painel e a estimação de associações condicionais, considerando pesos amostrais, controles e efeitos fixos. A fundamentação distingue exposição potencial, adoção efetiva da tecnologia e resultados observados no emprego. A análise deverá examinar a magnitude das associações e suas limitações, incluindo mudanças na coleta e na composição da amostra. Esta versão constitui um esqueleto provisório para o exercício ia-6.1 e não apresenta resultados empíricos novos. A contribuição pretendida consiste em organizar uma investigação reproduzível sobre o tema, sem atribuir causalidade automaticamente às diferenças observadas. As conclusões serão formuladas após a verificação dos resultados e a avaliação dos limites metodológicos pelo autor.
\end{resumo}